\section{Results and Discussion}
\label{sec:results_discussion}

\subsection{Primary End-to-End Comparison}

The primary comparison evaluated the flat four-class classifier against the
shared hierarchy under its deployable predicted-gate condition on the common
locked 3,668-image internal-test cohort. Table~\ref{tab:primary_results}
reports the frozen end-to-end results. Oracle routing is included only as a
diagnostic condition and is not deployable performance.

\begin{table}[t]
    \centering
    \caption{Frozen four-class internal-test performance. Oracle routing uses
    the ground-truth broad category and is diagnostic only.}
    \label{tab:primary_results}
    \footnotesize
    \setlength{\tabcolsep}{2.2pt}
    \begin{tabular}{lcccc}
        \toprule
        System & Acc. & BAcc. & M-F1 & W-F1 \\
        \midrule
        Flat four-class
        & 0.7421 & 0.6503 & \textbf{0.6192} & 0.7526 \\
        Shared hierarchy, predicted gate
        & 0.6876 & 0.6395 & 0.5686 & 0.7068 \\
        Shared hierarchy, oracle gate
        & \textbf{0.9338} & \textbf{0.7792} & \textbf{0.7770} & \textbf{0.9340} \\
        \bottomrule
    \end{tabular}
\end{table}

Under deployable predicted routing, the flat classifier achieved macro-F1
0.619222 and accuracy 0.742094, compared with 0.568591 and 0.687568 for the
shared hierarchy. The flat classifier therefore had the stronger observed
end-to-end performance under the frozen internal evaluation. The oracle row
must not be interpreted as a competing operational system because it uses the
ground-truth route.

\subsection{Paired Statistical Comparison}

The hierarchy-minus-flat macro-F1 difference was $-0.050632$. A 10,000-
replicate ground-truth-class-stratified paired bootstrap gave a percentile
95\% confidence interval of $[-0.075993,\,-0.024827]$. The corresponding
accuracy difference was $-0.054526$, with 95\% confidence interval
$[-0.068702,\,-0.040349]$. Both intervals excluded zero and favored the flat
classifier on the frozen internal test.

At the sample level, both systems were correct for 2,248 images and both were
wrong for 672. The flat classifier alone was correct for 474 images, whereas
the hierarchy alone was correct for 274. Exact two-sided McNemar testing gave
$p=2.49118\times10^{-13}$, supporting asymmetric paired correctness in favor
of the flat classifier. McNemar's test concerns paired correctness and the
accuracy comparison; it is not a significance test for macro-F1.

\subsection{Class-Specific Behaviour}

\begin{table}[t]
    \centering
    \caption{Per-class F1 on the locked internal test. Differences and
    confidence intervals are hierarchy minus flat.}
    \label{tab:per_class_comparison}
    \footnotesize
    \setlength{\tabcolsep}{2.8pt}
    \begin{tabular}{lrrrr}
        \toprule
        Class & Support & Hier. F1 & Flat F1 & $\Delta$ H--F \\
        \midrule
        Non-malignant & 2398 & 0.7786 & \textbf{0.8218} & -0.0432 \\
        Melanoma      &  678 & 0.5457 & \textbf{0.6094} & -0.0638 \\
        BCC           &  498 & 0.6592 & \textbf{0.6885} & -0.0293 \\
        SCC           &   94 & 0.2909 & 0.3571 & -0.0662 \\
        \bottomrule
    \end{tabular}
\end{table}

The flat classifier had higher F1 point estimates for all four endpoint
classes. The hierarchy-minus-flat 95\% confidence intervals were
$[-0.056035,-0.030510]$ for non-malignant, $[-0.088293,-0.039765]$ for
melanoma, and $[-0.056238,-0.002722]$ for BCC. For SCC, the point estimate
also favored the flat classifier, but the interval was
$[-0.150000,0.019586]$ and crossed zero. With only 94 SCC test images, the
evidence does not support a reliable SCC-specific superiority claim.

\subsection{Routing-Error Decomposition}

Oracle routing retained the trained shared downstream prediction heads but
replaced the learned broad routing decision with the ground-truth broad
category. This diagnostic increased four-class macro-F1 from 0.568591 under
predicted routing to 0.776976, producing a routing-associated gap of 0.208385.

\begin{table}[t]
    \centering
    \caption{Frozen routing decomposition for the shared hierarchy.}
    \label{tab:routing_results}
    \footnotesize
    \setlength{\tabcolsep}{3.2pt}
    \begin{tabular}{lr}
        \toprule
        Quantity & Result \\
        \midrule
        Predicted-gate macro-F1 & 0.568591 \\
        Oracle-gate macro-F1 & 0.776976 \\
        Routing-associated gap & \textbf{0.208385} \\
        \midrule
        Malignant blocked at Task~1 & 170/1,270 (13.39\%) \\
        Non-malignant sent to Task~2 & 761/2,398 (31.73\%) \\
        Wrong subtype after correct route & 215/1,100 (19.55\%) \\
        \bottomrule
    \end{tabular}
\end{table}

Of 1,270 malignant cases, 170 were blocked at Task~1 and therefore could not
receive a malignant subtype prediction. Of 2,398 non-malignant cases, 761
were incorrectly routed to Task~2, where only malignant outputs were
available. Among 1,100 correctly routed malignant cases, 215 received the
wrong subtype. These counts, together with the 0.208385 predicted-versus-
oracle macro-F1 gap, identify routing error as a major end-to-end bottleneck,
but they do not establish routing as the sole cause of hierarchical
underperformance.

Conditional subtype performance remained materially stronger than deployed
end-to-end hierarchical performance. On the true malignant subset, the shared
Task~2 head achieved macro-F1 0.702634 and accuracy 0.808661. The distinction
between conditional Task~2 capability and deployed four-class performance is
therefore important.

\subsection{Shared Versus Standalone Task Models}

\begin{table}[t]
    \centering
    \caption{Shared versus standalone task macro-F1. These comparisons are
    descriptive and do not establish causal negative transfer.}
    \label{tab:shared_standalone}
    \footnotesize
    \setlength{\tabcolsep}{3.3pt}
    \begin{tabular}{lrrr}
        \toprule
        Task & Shared & Standalone & Shared--Standalone \\
        \midrule
        Task~1 & 0.740624 & \textbf{0.774009} & -0.033385 \\
        Task~2 & 0.702634 & \textbf{0.724875} & -0.022241 \\
        Task~3 & \textbf{0.298710} & 0.275611 & +0.023099 \\
        \bottomrule
    \end{tabular}
\end{table}

Shared Task~1 and Task~2 underperformed their standalone counterparts in
macro-F1 by 0.033385 and 0.022241, respectively. This pattern is consistent
with possible negative transfer or insufficient task-specific specialization,
but the descriptive comparisons do not establish causation and were not
supported by task-specific paired confidence intervals.

Task~3 showed the opposite macro-F1 direction, with shared macro-F1 0.298710
versus 0.275611 for the standalone model. However, the Task~3 test subset had
only 127 cases, with T2, T3, and T4 supports of 7, 2, and 1. These rare
category estimates are therefore highly unstable and do not support strong
T-category conclusions.

\subsection{Validation-to-Test Directional Consistency}

The primary comparison had the same direction on validation and on the final
internal test. The hierarchy-minus-flat macro-F1 difference was $-0.070285$
on validation and $-0.050632$ on the internal test; the corresponding
accuracy differences were $-0.063795$ and $-0.054526$. Routing loss was
0.192480 on validation and 0.208385 on the internal test. This is evidence of
within-study directional consistency only. The internal-test results were not
used for additional model selection, threshold tuning, calibration, or
retraining.

\subsection{Limitations}

The final evidence comes from one internal ISIC-derived test partition rather
than an external, prospective, multi-centre, or clinical evaluation. The split
controlled available lesion identifiers and exact duplicates, but fully
patient-independent separation could not be guaranteed because patient
identifiers were unavailable and 2,084 images lacked lesion identifiers.

Minority-class uncertainty remains substantial. SCC support was only 94 and
its paired class-specific interval crossed zero. Task~3 was still more sparse,
with T2/T3/T4 supports of 7/2/1. The bootstrap intervals quantify uncertainty
under resampling of the observed test cases but do not capture retraining,
seed, acquisition-site, label-noise, or distribution-shift uncertainty.

No external evaluation, calibration study, explainability experiment, or
skin-tone fairness analysis was completed. The evidence therefore supports a
controlled internal comparative and failure-analysis study, not clinical
validation, external generalisation, or deployment readiness.
